\documentclass[11pt,a4paper]{article}

\usepackage[utf8]{inputenc}
\usepackage[T5]{fontenc}
\usepackage[vietnamese]{babel}
\usepackage[margin=2cm]{geometry}
\usepackage{fontawesome5}
\usepackage{xcolor}
\usepackage{titlesec}
\usepackage{enumitem}
\usepackage{hyperref}
\usepackage{parskip}

% Define colors
\definecolor{primary}{RGB}{204,85,0} % Burnt Orange
\definecolor{secondary}{RGB}{60,60,60} % Dark Gray

% Font settings (optional, stick to default for compatibility or use helvet if desired)
% \usepackage{helvet}
% \renewcommand{\familydefault}{\sfdefault}

% Custom section formatting
\titleformat{\section}
{\Large\bfseries\color{primary}}
{}
{0em}
{}[\titlerule]

\titlespacing{\section}{0pt}{12pt}{8pt}

\titleformat{\subsection}
{\large\bfseries\color{secondary}}
{}
{0em}
{}

\titlespacing{\subsection}{0pt}{8pt}{4pt}

% Custom itemize
\setlist[itemize]{leftmargin=*}

\begin{document}

% HEADER
\begin{center}
    {\Huge\bfseries\color{primary} NGUYỄN VY THÔNG}\\[5pt]
    {\Large Giáo viên Toán | Giáo viên Toán có chuyên môn về Công nghệ Giáo dục AI | Cử nhân Toán học Tài năng}\\[10pt]
    \faMapMarker*~TP. Hồ Chí Minh, Việt Nam \quad
    \faPhone~(+84) 394 720 727 \quad
    \faEnvelope~\href{mailto:nvythong@gmail.com}{nvythong@gmail.com}\\[5pt]
    \faLink~\href{https://drive.google.com/drive/folders/1tL5VO2r13Qb-Fy3bYrKF97-6ZKrWLgaq?usp=drive_link}{Kho Minh Chứng Online: Link Google Drive}
\end{center}

% SUMMARY
\section{\faClipboardList~TÓM TẮT MỤC TIÊU}
Cử nhân Toán học với nền tảng chuyên sâu về Lý thuyết Tối ưu hóa và Thuật toán (GPA 9.26/10), kết hợp với kinh nghiệm thực tiễn trong giảng dạy và nghiên cứu ứng dụng Trí tuệ nhân tạo (AI). 
Là một nhà giáo dục tâm huyết, tôi hướng tới việc đổi mới phương pháp giảng dạy toán học thông qua công nghệ, đặc biệt là tích hợp tư duy máy tính và AI vào giáo dục phổ thông. Sở hữu khả năng tự học mạnh mẽ, tư duy phản biện và niềm đam mê với khoa học nhận thức.

% EDUCATION
\section{\faGraduationCap~HỌC VẤN}

\textbf{Cử nhân Toán học} \hfill (9/2016 - 10/2020)\\
\textit{Trường Đại học Khoa học Tự nhiên, ĐHQG-HCM}
\begin{itemize}
    \item \textbf{Điểm trung bình (GPA):} 9.26/10 (Xuất sắc)
    \item \textbf{Chuyên ngành:} Tối ưu hóa và Lý thuyết hệ thống.
    \item \textbf{Luận văn tốt nghiệp:} ``Tập tiếp xúc bậc cao và đạo hàm bậc cao của ánh xạ giá trị tập hợp''.
    \item \textbf{Thành tích nổi bật:}
    \begin{itemize}
        \item Học bổng Chương trình trọng điểm Quốc gia phát triển Toán học (VIASM) (Từ năm 2017 đến năm 2020).
        \item Học bổng Khuyến khích học tập của ĐHKHTN (3 năm liên tiếp).
        \item Lớp trưởng Xuất sắc của năm (2016 - 2018).
    \end{itemize}
\end{itemize}

\textbf{Cao đẳng Marketing} \hfill (9/2011 - 9/2014)\\
\textit{Trường Cao đẳng Kinh tế Đối ngoại}
\begin{itemize}
    \item \textbf{GPA:} 6.85/10
\end{itemize}

% EXPERIENCE
\section{\faBriefcase~KINH NGHIỆM LÀM VIỆC}

\textbf{Giáo viên Toán thỉnh giảng} \hfill (3/2023 - Hiện tại)\\
\textit{4 trường tại khu vực huyện Củ Chi (cũ), TP.HCM}

\textit{Sinh ra, lớn lên và làm việc tại khu vực huyện Củ Chi (cũ), tôi đã và đang đảm nhận vai trò giảng dạy tại 4 trường với các nhiệm vụ đa dạng:}
\begin{itemize}
    \item \textbf{THCS Tân Phú Trung:} Bồi dưỡng Học sinh giỏi Toán \& Luyện thi Tuyển sinh lớp 10
    \item \textbf{THCS Tân Thạnh Tây \& THCS Tân Thông Hội:} Giảng dạy môn Toán theo chương trình GDPT
    \item \textbf{THPT Chiến Thắng:} Phụ trách Giảng dạy Chuyển đổi số \& AI (khối 10-12)
    \item \textbf{Thành tích:}
    \begin{itemize}
        \item Tham gia công tác \textbf{Bồi dưỡng Học sinh giỏi} môn Toán (Năm học 2023 - 2024).
        \item Đổi mới phương pháp dạy học, tích hợp công nghệ AI vào bài giảng.
    \end{itemize}
\end{itemize} 

\textbf{Trưởng phòng Học thuật - Chuyên viên Nghiên cứu \& Phát triển (R\&D)} \hfill (10/2021 - 7/2022)\\
\textit{Trung tâm AI (AIC), Khu Công nghệ Phần mềm ĐHQG-HCM (ITP)}
\begin{itemize}
    \item Nghiên cứu và phát triển các chương trình đào tạo về Trí tuệ nhân tạo (AI) dành cho học sinh và sinh viên.
    \item Quản lý chất lượng học thuật và đội ngũ giảng viên/trợ giảng.
    \item \textbf{Thành tích:} Đạt danh hiệu \textbf{Lao động tiên tiến} tại Khu Công nghệ Phần mềm ĐHQG-HCM.
\end{itemize}

\textbf{Trợ giảng (Teaching Assistant)} \hfill (10/2020 - 8/2021)\\
\textit{Trường Đại học Khoa học Tự nhiên, ĐHQG-HCM}
\begin{itemize}
    \item Hỗ trợ giảng dạy các học phần: Vi Tích Phân II , Quy hoạch tuyến tính, Các thuật toán tối ưu.
    \item Hướng dẫn sinh viên giải bài tập, thực hành lập trình và giải đáp thắc mắc chuyên môn.
\end{itemize}

\textbf{Giáo viên Toán} \hfill (6/2015 - 1/2016)\\
\textit{Trung tâm Mathnasium - Điện Biên Phủ}
\begin{itemize}
    \item Giảng dạy toán tư duy cho học sinh tiểu học (phương pháp Mathnasium).
    \item Hỗ trợ công tác quản lý trung tâm và tư vấn phụ huynh.
\end{itemize}

% PROJECTS
\section{\faRocket~DỰ ÁN NỔI BẬT}

\textbf{Hệ thống E-Learning AI cho Giảng dạy Toán THCS} \hfill (2023 - Nay)
\begin{itemize}
    \item \textbf{Tích hợp AI Agent:} Sử dụng Gemini/NotebookLM làm trợ lý ảo (TutorBot) hỗ trợ giải đáp toán học 24/7 cho học sinh.
    \item \textbf{Tự động hóa:} Phát triển công cụ sinh đề kiểm tra tự động bằng LaTeX/TikZ từ ngân hàng câu hỏi số.
    \item \textbf{Hiệu quả:} Giảm 50\% thời gian soạn bài, tăng mức độ tương tác học sinh thông qua học tập cá nhân hóa.
\end{itemize}

\textbf{Chương trình Đào tạo AI và Robotic (Tại AIC Center)} \hfill (2021 - 2022)
\begin{itemize}
    \item Thiết kế khung chương trình giáo dục AI cho đối tượng K-12 và sinh viên không chuyên.
    \item Nghiên cứu ứng dụng Robotic trong Giáo dục Nhận thức về Trí Thông Minh Nhân tạo ở cấp Tiểu học.
\end{itemize}

% CERTIFICATES
\section{\faCertificate~CHỨNG CHỈ \& ĐÀO TẠO}

\subsection*{Chứng chỉ Sư phạm}
\begin{itemize}
    \item \textbf{Chứng chỉ Bồi dưỡng nghiệp vụ sư phạm Giáo viên Toán THCS} (ĐH Sư phạm Hà Nội 2, 2024).
\end{itemize}

\subsection*{Chứng chỉ Quốc tế \& Chuyên môn}

\textbf{1. Sư phạm AI \& Lãnh đạo chuyển đổi số (AI Pedagogy \& Leadership)}
\begin{itemize}
    \item \textbf{Navigating Generative AI for Leaders} (Vanderbilt Univ.): Kỹ năng ra quyết định và triển khai chiến lược AI.
    \item \textbf{Prompt Engineering for Educators} (Vanderbilt Univ.): Kỹ thuật soạn giảng và tương tác AI nâng cao.
    \item \textbf{Google Gemini Certified Educator} (Google): Chứng nhận năng lực ứng dụng Gemini trong giáo dục.
    \item \textbf{Teaching AI Fluency Framework} (Anthropic): Khung năng lực đào tạo AI cho người học.
    \item \textbf{Generative AI \& ChatGPT for K-12 Educators}: Ứng dụng AI thực tiễn trong giáo dục phổ thông.
\end{itemize}

\textbf{2. Kỹ thuật AI chuyên sâu (Technical Depth)}
\begin{itemize}
    \item \textbf{Neural Networks and Deep Learning} (DeepLearning.AI).
    \item \textbf{Statistical Learning} (Vietnam Institute of Mathematics - VIASM).
    \item \textbf{Data Science \& Machine Learning Specializations} (Stanford Online / DeepLearning.AI).
\end{itemize}

\textbf{3. Ngoại ngữ}
\begin{itemize}
    \item \textbf{TOEIC:} 865/990 (2023) - Sử dụng thành thạo tiếng Anh chuyên ngành.
\end{itemize}

\subsection*{Khóa học Kỹ năng Bổ trợ}
\begin{itemize}
    \item \textbf{Sư phạm \& Kỹ năng mềm:} Online Education (Macquarie Univ.), Create Video/Audio for Online Learning, Philosophy of Cognitive Sciences.
    \item \textbf{Hoạt động học thuật:} Trường hè Toán học 2020 (Viện Toán học), Trường hè Toán học 2019 (VIASM), Qiskit Global Summer School 2020 (IBM Quantum).
\end{itemize}

% SKILLS
\section{\faTools~KỸ NĂNG \& NĂNG LỰC CỐT LÕI}

\textbf{1. Trí tuệ nhân tạo \& Khoa học Dữ liệu (AI \& Data Science)}
\begin{itemize}
    \item \textbf{Generative AI Mastery:} Thành thạo Prompt Engineering \& Context Engineering (Claude/Gemini/GPT-5/Qwen/DeepSeek) để xây dựng trợ lý ảo và Agentic workflows.
    \item \textbf{Technical Stack:} Python, Julia, MATLAB.
    \item \textbf{Data Analysis:} Phân tích dữ liệu giáo dục, Tối ưu hóa mô hình học tập cá nhân hóa.
\end{itemize}

\textbf{2. Công nghệ Giáo dục \& Sư phạm số (EdTech \& Digital Pedagogy)}
\begin{itemize}
    \item \textbf{AI-Integrated Curriculum:} Thiết kế chương trình học tích hợp AI, sử dụng các Framework sư phạm hiện đại (Scaffolding, Bloom's Taxonomy).
    \item \textbf{Digital Content Creation:} Soạn thảo tài liệu Toán học chất lượng cao với LaTeX/TikZ, thiết kế E-Learning tương tác.
    \item \textbf{Assessment Tech:} Xây dựng hệ thống đánh giá tự động và phản hồi tức thì cho học sinh.
\end{itemize}

\textbf{3. Tư duy Toán học \& Nghiên cứu (Mathematical \& Research Mindset)}
\begin{itemize}
    \item \textbf{Advanced Math:} Tối ưu hóa, Lý thuyết hệ thống, Mô hình hóa toán học.
    \item \textbf{Research:} Khả năng đọc hiểu và chuyển ngữ các tài liệu/nghiên cứu AI tiên tiến nhất từ tiếng Anh.
\end{itemize}

% AWARDS
\section{\faTrophy~GIẢI THƯỞNG \& VINH DANH KHÁC}
\begin{itemize}
    \item Danh hiệu Lao động tiên tiến - Khu Công nghệ Phần mềm ĐHQG-HCM.
    \item Các học bổng academic uy tín từ VIASM và ĐHKHTN.
\end{itemize}

\end{document}
